\subsection{伴随素理想的定义}

定义 1. 设 M 是环 A 上的模。如果存在 \(x \in M\)，使得 p 等于 x 的湮灭子，则称素理想 p 与 M 相伴。与 M 相伴的素理想所成的集合记为 \(\operatorname{Ass}_{\mathbf{A}}(\mathbf{M})\)，或简称为 \(\operatorname{Ass}(\mathbf{M})\)。

* 例。设 \(a\) 是多项式环 \(A = \mathbf{C}[X_1, \ldots, X_r]\) 中的一个理想，\(V\) 是相应的仿射代数簇，\(V_1, \ldots, V_p\) 是 \(V\) 的不可约分支。若取 \(M\) 为在 \(A/a\) 上正则的函数环 \(V\)，则与 \(M\) 相伴的素理想集合由 \(V_1, \ldots, V_p\) 的理想以及一般还包括其他素理想组成，后者中的每一个都包含某个 \(V_1\) 的理想。

由于 0 的湮灭子是 A，湮灭子为素理想的元素 \(x \in \mathbf{M}\) 必然满足 \(\neq 0\)。称素理想 \(p\) 与 M 相伴，等价于说 M 含有一个同构于 A/p 的子模（即对于湮灭子为 p 的所有 \(x \in M\)，子模 Ax）。

如果 A-模 M 是子模族 \((M_{t})_{t\in I}\) 的并，则显然

\[
\operatorname{Ass} (\mathbf {M}) = \bigcup_ {\iota \in I} \operatorname{Ass} \left(\mathbf {M} _ {\iota}\right).\tag{1}
\]

命题 1. 对于每个素理想 \(\mathfrak{p}\)（在环 \(\mathbf{A}\) 中）以及 \(\mathbf{M} \neq 0\) 的每个子模 \(\mathbf{A} / \mathfrak{p}\)，都有 \(\operatorname{Ass}(\mathbf{M}) = \{\mathfrak{p}\}\)。

由于环 A/p 是整环，A/p 中任一元素 \(\neq0\) 的湮灭子都是 p。

命题 2. 设 M 是环 A 上的模。由 M 中所有元素 \(\neq0\) 所对应的 A 的理想 Ann(x) 构成的集合的每个极大元都属于 Ass(M)。

令 \(\mathfrak{a} = \operatorname{Ann}(x)\) ( \(x \in \mathbf{M}, x \neq 0\) ) 为这样的极大元；只需证明 \(\mathfrak{a}\) 是素理想。由于 \(x \neq 0\)，\(\mathfrak{a} \neq \mathbf{A}\)。设 \(b, c\) 是 \(\mathbf{A}\) 中满足 \(bc \in \mathfrak{a}\) 且 \(c \notin \mathfrak{a}\) 的元素。则 \(cx \neq 0\)，\(b \in \operatorname{Ann}(cx)\) 且 \(\mathfrak{a} \subset \operatorname{Ann}(cx)\)。由于 \(\mathfrak{a}\) 是极大的，\(\operatorname{Ann}(cx) = \mathfrak{a}\)，从而 \(b \in \mathfrak{a}\)，故 \(\mathfrak{a}\) 是素理想。

推论 1. 设 M 是诺特环 A 上的模。则条件 \(M \neq \{0\}\) 等价于 \(\operatorname{Ass}(M) \neq \varnothing\)。

如果 \(M = \{0\}\)，显然 \(\text{Ass}(M)\) 为空（对 A 不需作任何假设）。如果 \(M \neq \{0\}\)，则形如 \(\text{Ann}(x)\)（其中 \(x \in M\) 且 \(x \neq 0\)）的理想所成的集合非空，且其中的理想都满足 \(\neq A\)；由于 A 是诺特环，该集合有极大元；于是只需应用命题 2。

推论 2. 设 A 是诺特环，M 是 A-模，a 是 A 的元素。M 上比为 a 的伸缩映射为单射的充要条件是：a 不属于任何与 M 相伴的素理想。

如果 \(a\) 属于素理想 \(\mathfrak{p} \in \operatorname{Ass}(\mathbf{M})\)，则 \(\mathfrak{p} = \operatorname{Ann}(x)\)，其中 \(x \in \mathbf{M}\)，\(x \neq 0\)；于是 \(ax = 0\)，乘为 \(a\) 的伸缩映射不是单射。反之，若 \(ax = 0\) 对某个 \(x \in \mathbf{M}\) 且 \(x \neq 0\) 成立，则 \(Ax \neq \{0\}\)，从而 \(\operatorname{Ass}(x) \neq \varnothing\)（推论 1）。令 \(\mathfrak{p} \in \operatorname{Ass}(\mathbf{A}x)\)；则显然 \(\mathfrak{p} \in \operatorname{Ass}(\mathbf{M})\)，且 \(\mathfrak{p} = \operatorname{Ann}(bx)\)，其中 \(b \in \mathbf{A}\)；故 \(a \in \mathfrak{p}\)，因为 \(abx = 0\)。

推论 3. 诺特环 A 中的零因子集合是理想 \(\mathfrak{p} \in \operatorname{Ass}(A)\) 的并。

命题 3. 设 A 是一个环，M 是一个 A-模，N 是 N 的子模。则

\[
\operatorname{Ass} (\mathbf {N}) \subset \operatorname{Ass} (\mathbf {M}) \subset \operatorname{Ass} (\mathbf {N}) \cup \operatorname{Ass} (\mathbf {M} / \mathbf {N}).\tag{2}
\]

包含关系 \(\operatorname{Ass}(\mathbf{N}) \subset \operatorname{Ass}(\mathbf{M})\) 显然。设 \(\mathfrak{p} \in \operatorname{Ass}(\mathbf{M})\)，\(\mathbf{E}\) 是 \(\mathbf{M}\) 中同构于 \(\mathbf{A} / \mathfrak{p}\) 的子模，且 \(\mathbf{F} = \mathbf{E} \cap \mathbf{N}\)。若 \(\mathbf{F} = \{0\}\)，则 \(\mathbf{E}\) 同构于 M/N 的一个子模，从而 \(\mathfrak{p} \in \operatorname{Ass}(M/N)\)。若 \(F \neq \{0\}\)，则 F 中每个元素 \(\neq 0\) 的湮灭子都是 p（命题 1），因而 \(\mathfrak{p} \in \operatorname{Ass}(F) \subset \operatorname{Ass}(N)\)。

推论 1. 如果 A-模 M 是子模族 \((\mathbf{M}_t)_{t \in I}\) 的直和，则 \(\operatorname{Ass}(\mathbf{M}) = \bigcup_{t \in I} \operatorname{Ass}(\mathbf{M}_t)\)。

由 (1) 可将其归结为 I 是有限集的情形，再通过对\(\operatorname{Card}(I) = 2\)归纳将其归结为\(\operatorname{Card}(I)\)的情形。令\(I = \{i, j\}\)，其中\(i \neq j\)；由于\(M/M_{i}\)同构于\(M_{j}\)，\(\operatorname{Ass}(M) \subset \operatorname{Ass}(M_{i}) \cup \operatorname{Ass}(M_{j})\)（命题 3）；此外，\(\operatorname{Ass}(M_{i})\)和\(\operatorname{Ass}(M_{j})\)都包含于\(\operatorname{Ass}(M)\)（命题 3），结论得证。

推论 2. 设 M 是 A-模，\((\mathbf{Q}_i)_{i\in I}\) 是 M 的有限个子模组成的族。若 \(\bigcap_{i\in I} \mathbf{Q}_i = \{0\}\)，则

\[
\operatorname{Ass} (\mathbf {M}) \subset \bigcup_ {i \in I} \operatorname{Ass} (\mathbf {M} / \mathbf {Q} _ {i}).
\]

典范映射 \(\mathbf{M} \to \bigoplus_{i \in I} (\mathbf{M} / \mathbf{Q}_i)\) 是单射；于是只需应用命题 3 及其推论 1。

命题 4. 设 M 是 A-模，\(\Phi\) 是 \(\operatorname{Ass}(\mathbf{M})\) 的子集。则存在 M 的一个子模 N，使得 \(\operatorname{Ass}(\mathbf{N}) = \operatorname{Ass}(\mathbf{M}) - \Phi\) 且 \(\operatorname{Ass}(\mathbf{M}/\mathbf{N}) = \Phi\)。

令\(\mathfrak{E}\)为 M 的子模\(P\)的集合，满足\(M\)。公式 (1) 表明，按包含关系排序的集合\(\operatorname{Ass}(P) \subset \operatorname{Ass}(M) - \Phi\)是归纳的；此外，\(\mathfrak{E}\)，因而\(\{0\} \in \mathfrak{E}\)。令\(\mathfrak{E} \neq \varnothing\)为\(N\)的极大元。则\(\mathfrak{E}\)。我们将看到\(\operatorname{Ass}(N) \subset \operatorname{Ass}(M) - \Phi\)，命题 3 将完成证明。设\(\operatorname{Ass}(M/N) \subset \Phi\)；则\(p \in \operatorname{Ass}(M/N)\)含有一个子模\(M/N\)同构于\(F/N\)。由命题 1 和命题 3，\(A/p\)。由于\(\operatorname{Ass}(F) \subset \operatorname{Ass}(N) \cup \{p\}\)是\(N\)中的极大元，\(\mathfrak{E}\)不属于\(F \notin \mathfrak{E}\)，从而\(p \in \Phi\)。

